\chapter{The Sketcher: Constraint-Based Profiles}
\label{ch:sketcher}

Almost every solid built for analysis begins as a 2D \emph{sketch}: a
planar collection of geometry (lines, arcs, circles, B-splines)
governed by \emph{constraints}. The Sketcher is a
\textbf{2D geometric constraint solver}, and understanding it as such ---
rather than as a drawing tool --- is the key to using it well.

\section{Degrees of freedom and the solver}

A sketch's geometry carries \emph{degrees of freedom} (DOF): a free point
has 2 (its $x,y$), a free line segment has 4, and so on. Constraints
remove DOF. The Sketcher continuously reports the remaining DOF in the
solver-messages area:

\begin{itemize}
  \item \textbf{Under-constrained} ($\text{DOF}>0$): geometry can still
        move; dragging changes shape.
  \item \textbf{Fully constrained} ($\text{DOF}=0$): geometry is fixed
        and the sketch turns green. This is the target state for a
        production model.
  \item \textbf{Over-constrained / conflicting}: redundant or
        contradictory constraints; the solver flags them and refuses to
        solve. Redundant constraints should be removed, not left.
\end{itemize}

Mathematically the solver assembles the constraints as a system of
(generally nonlinear) equations $\vec{c}(\vec{q})=\vec 0$ in the geometry
parameters $\vec q$ and solves it by a Newton-type iteration; the DOF
count is $\dim\vec q$ minus the rank of the constraint Jacobian.

\section{Two kinds of constraint}

\begin{description}
  \item[Geometric constraints] impose relationships without numbers:
        coincident, horizontal, vertical, parallel, perpendicular,
        tangent, equal, symmetric, point-on-object.
  \item[Dimensional constraints] impose measured values: distance,
        horizontal/vertical distance, radius, diameter, angle. These are
        the \emph{parameters} that make the model editable, and they can
        be named and driven by expressions or a spreadsheet.
\end{description}

\begin{fctip}[Constrain like an engineer]
Prefer a small set of dimensional constraints tied to design intent
(``bracket width'', ``hole pitch'') over many arbitrary ones. Name the
important dimensions; you can then reference them in expressions
(\texttt{Sketch.Constraints.width}) and sweep them for a parametric
study --- exactly what a design-of-experiments FEA campaign needs.
\end{fctip}

\section{Workflow}

The typical loop is: pick a plane (or a planar face) $\rightarrow$ enter
the sketch $\rightarrow$ draw rough geometry $\rightarrow$ add geometric
constraints to capture intent $\rightarrow$ add dimensional constraints
to fix size $\rightarrow$ reach $\text{DOF}=0$ $\rightarrow$ close the
sketch. The finished sketch becomes the profile for a solid feature
(\cref{ch:part-partdesign}).

\section{External geometry and construction geometry}

\emph{Construction geometry} (toggle mode) participates in constraints
but is not part of the output profile --- useful for reference lines and
symmetry axes. \emph{External geometry} imports edges from existing
solids into the sketch as references, letting a new feature key off
previously built geometry while preserving the parametric link.
